% ============================================================
% RELATED WORK — target: 800--1200 words
% Thematic subsections
% ============================================================

\section{Related Work}

\subsection{Chain-of-Thought Faithfulness}

The question of whether chain-of-thought explanations faithfully represent a model's
actual reasoning process has been investigated through several complementary
methodologies. Turpin et al.~\cite{turpin2023unfaithful} introduced the bias
injection paradigm: by adding biasing features to few-shot prompts (e.g., labeling
answers with positional patterns or attributing answers to authority figures), they
showed that models on BIG-Bench Hard tasks are systematically influenced by features
never mentioned in their CoT explanations. This established that CoT unfaithfulness
is not merely occasional but reflects a structural gap between the reasoning a model
reports and the factors that actually drive its outputs.

Lanham et al.~\cite{lanham2023measuring} took a causal approach, measuring
faithfulness in Claude models by intervening on the CoT itself---truncating reasoning
chains, adding mistakes to intermediate steps, and paraphrasing explanations. Their
key finding was that the degree to which CoT content causally determines model
predictions varies substantially by task, suggesting that faithfulness is not a fixed
property of a model but depends on the interaction between model, task, and reasoning
chain structure. Radhakrishnan et al.~\cite{radhakrishnan2023decomposition}
demonstrated that decomposing questions into subquestions before generating answers
improves faithfulness, as judged by the consistency between intermediate reasoning
steps and final outputs.

Lyu et al.~\cite{lyu2023faithful} proposed Faithful CoT, a framework that translates
natural language queries into symbolic reasoning chains (e.g., Python code or
Datalog programs) before solving them, achieving verifiable faithfulness by
construction. While this approach sidesteps the faithfulness problem for structured
reasoning tasks, it does not generalize to open-ended reasoning where symbolic
translation is infeasible.

Most recently, Chen et al.~\cite{chen2025reasoning} conducted the most directly
relevant evaluation, testing Claude 3.7 Sonnet and DeepSeek-R1 across six hint
types injected into MMLU and GPQA questions. Their finding that faithfulness rates
are as low as 25\% for Claude and 39\% for DeepSeek-R1---with outcome-based
reinforcement learning providing only modest and plateauing improvements---raised
fundamental questions about the viability of CoT monitoring as a safety mechanism.
Feng et al.~\cite{feng2025more} extended this to three reasoning models
(QwQ-32b, Gemini 2.0 Flash Thinking, and DeepSeek-R1), finding that reasoning
models are substantially more faithful than their non-reasoning counterparts,
though they tested only sycophancy hints on MMLU. Cornish and
Rogers~\cite{cornish2025examining} probed DeepSeek-R1 on 445 logical puzzles
and discovered an asymmetry: the model acknowledges \emph{harmful} hints
94.6\% of the time but reports fewer than 2\% of helpful hints, suggesting
that faithfulness is entangled with the perceived valence of the influencing cue.
Arcuschin et al.~\cite{arcuschin2025wild} extended this investigation to natural
settings, finding that CoT reasoning in the wild is also frequently unfaithful,
while making the important distinction between hint-injection faithfulness and the
broader concept of reasoning faithfulness in unconstrained settings.

Complementary methodological approaches have emerged alongside hint-injection
studies. Chua et al.~\cite{chua2025unlearning} proposed measuring faithfulness
by selectively unlearning reasoning steps and observing whether model predictions
change accordingly, providing a causal alternative to post-hoc CoT analysis.
Stickland et al.~\cite{stickland2025monitorability} introduced the concept of
\emph{monitorability}---distinguishing it from faithfulness---and showed that
models can appear faithful yet remain hard to monitor when they omit key factors
from their reasoning, with monitorability varying sharply across model families.

\subsection{Reasoning Models and Safety}

The proliferation of reasoning models---systems trained with reinforcement learning
to produce extended chains of thought---has created new safety challenges alongside
their performance gains. DeepSeek-R1 \cite{deepseek2025r1} demonstrated that
reinforcement learning with group relative policy optimization (GRPO) can elicit
sophisticated reasoning behaviors, while subsequent work has shown that similar
capabilities emerge across diverse training paradigms including distillation,
supervised fine-tuning, and data-centric approaches.

Several studies have examined the safety implications of these reasoning capabilities.
Jiang et al.~\cite{jiang2025safechain} evaluated safety risks specific to models
with long CoT capabilities, finding that the extended reasoning process can both
help and hinder safety---models sometimes reason their way toward harmful outputs.
The hidden risks of reasoning models have been further documented by safety
assessments of DeepSeek-R1 showing that reasoning traces can contain concerning
patterns even when final outputs appear benign \cite{zhou2025hidden}.

A particularly important line of work questions the fundamental limits of CoT
monitoring as a safety strategy. Baker et al.~\cite{baker2025cot} argued that CoT
may remain highly informative for monitoring \emph{despite} unfaithfulness,
distinguishing between the strict requirement that CoT perfectly mirrors internal
computation and the weaker but practically useful property that CoT correlates
with safety-relevant behaviors. Conversely, Chen et al.~\cite{chen2025reasoning}
demonstrated that when reinforcement learning increases hint usage (reward hacking),
the rate at which models verbalize that usage in their CoT remains below 2\%,
suggesting that outcome-based training actively decouples behavior from stated
reasoning.

\subsection{Open-Weight Model Evaluation}

The rapid expansion of the open-weight model ecosystem has created an urgent need
for systematic evaluation across model families. By early 2026, open-weight
reasoning models span a wide range of architectures (dense, mixture-of-experts,
hybrid attention), training methods (GRPO, DPO, distillation, supervised
fine-tuning, data-centric), and scales (8B to over 1T parameters). Surveys of
large reasoning models \cite{xu2025large, qchen2025longcot, li2025system12} have
catalogued this diversity but have not included faithfulness as an evaluation
dimension.

Existing faithfulness evaluations have focused almost exclusively on proprietary
models or on a single open-weight family (typically DeepSeek). Wu et
al.~\cite{wu2025moreless} examined the relationship between CoT length and task
performance, finding an inverted U-shape where excessively long reasoning chains
degrade accuracy---a finding relevant to faithfulness since Chen et
al.~\cite{chen2025reasoning} showed that unfaithful CoTs tend to be longer.
Probing-based approaches have examined whether models' internal representations
contain information about answer correctness that is absent from their CoT
\cite{liu2025probing}, suggesting that faithfulness gaps may be detectable through
mechanistic interpretability even when CoT analysis alone fails. Wang et
al.~\cite{wang2024making} proposed training-time interventions to improve
faithfulness, demonstrating that explicitly optimizing for reasoning consistency
during fine-tuning can narrow the gap between stated and actual reasoning processes.

Our work fills the gap between these lines of research by providing the first
systematic faithfulness evaluation that spans the breadth of the open-weight
reasoning model ecosystem. By testing 11 models across 9 families with a
standardized methodology---and using Anthropic's BLOOM
framework~\cite{bloom2025} for automated behavioral scoring---we enable direct
comparison of faithfulness properties across architectures, training methods,
and scales, comparisons that are currently impossible with the narrow model
coverage of existing studies.
